\begin{abstract}
本文作者所倡导的大动量有效理论（Large-Momentum Effective Theory，简称 LaMET）为在欧氏格点
QCD 理论中模拟部分子物理提供了一条直接途径。近来，文献中对这一理论表现出浓厚兴趣，其中不乏
对其正确性与有效性的质疑。这里我们提供若干讨论，旨在进一步阐明该方法。特别地，我们解释了
为什么在格点 QCD 计算部分子分布矩时它不具有通常的幂次发散问题。LaMET 方法中唯一的幂次发散来自
威尔逊线的自能，而它可以被恰当地因子化出去。我们证明：虽然 Ioffe 时间分布为从同一组格点观测量
中提取部分子分布提供了另一种方式，但要获得精度计算，它同样需要与 LaMET 相同的大动量（或短距离）
极限。在对误差作恰当量化之后，两种提取方法应与同一批格点数据相比较。
\end{abstract}

\section{引言}
\label{intro}

费曼发明的部分子模型一直是描述高能强相互作用物理最有用的语言之一~\cite{Feynman:1969ej}。
部分子是强子以光速运动时才存在的理想化对象。因此用现代语言来说，它们是某个有效理论中的有效
自由度，例如软共线有效理论~\cite{Bauer:2000yr}。强子的部分子性质，如部分子分布、部分子波函数
或振幅、部分子关联等，都是低能性质，需要在强耦合区域求解量子色动力学（QCD）。

当在强子静止系中表述时，部分子物理表示 QCD 夸克或胶子沿光锥方向（$\xi^-=(t-z)/\sqrt{2}$）
的时空关联，它显含时间依赖~\cite{Sterman:1994ce}。（我们用 $\xi^\mu = (t, x, y, z)$
$(\mu=0,1,2,3)$ 表示时空坐标，尽管 $x$、$y$、$z$ 也常被用来表示动量分数。）在物理上，
含时关联自成一类，因为时间演化需要存在依赖相互作用的哈密顿量 $H$，这类可观测量被称为
``动力学的''。

动力学关联难以计算至少有两个不同的原因：1) 若插入一组中间哈密顿量本征态来简化演化，就需要知道
所有中间激发态的波函数。因此，仅仅知道初态（通常是基态）波函数不足以计算含时关联。为逼近基态
设计的方法未必适用于激发态，因而动力学关联的计算通常不如不含时的或静态的可观测量可靠——对后者
而言，初态波函数就足够了。2) 由于时间依赖涉及幺正演化算符 $e^{-iHt}$，它包含变号的相位，这对
数值蒙特卡罗模拟是个难题。事实上，大多数为计算静态性质设计的蒙特卡罗方法都无法直接用于动力学
关联。有人尝试改为计算虚时间或欧氏演化 $e^{-H\tau}$，再把结果解析延拓回实时间~\cite{Nakahara:1999vy,Asakawa:2000tr}。
然而这些方法很难可靠，因为其中的解析函数是未知的。另一种技巧是对时间演化算符作泰勒展开，通过
无穷级数之和获得时间依赖。由于高阶矩的计算精度很差，这种做法同样非常受限~\cite{Hagler:2009ni}。
基于上述原因，含时关联通常被称为``闵氏的''，意即时间关联与空间关联有本质区别。另一方面，
不含时关联可以称为``欧氏的''，因为它们可以方便地用蒙特卡罗模拟计算，只是可能存在额外的符号问题。

对于闵氏可观测量，传统做法是采用解析近似。例如，所谓的施温格--戴森方法通过动力学方程求解时间
依赖~\cite{Roberts:1994dr}。对于部分子物理，在所谓的光前量子化方法中，坐标被重新定义使新时间
成为空间与旧时间的组合，光锥关联于是变成不含时的~\cite{Burkardt:1995ct}。然而所得理论无法方便地
用蒙特卡罗方法求解，因为度规不是正定的。

所幸存在一些``动力学的''问题，它们本质上并非``闵氏的''——只要表述得当，就能用蒙特卡罗方法求解。
例如散射相移通常涉及复的 S 矩阵元，无法直接用蒙特卡罗方法计算；但可以把理论放进一个盒子里，
计算束缚态能级随边界条件的变化，并从束缚态能级提取相移~\cite{Luscher:1986pf,Luscher:1990ck,Luscher:1990ux,Briceno:2017max}。
第二个例子是某些闵氏关联函数可以直接解析延拓到欧氏区域而不遇到任何``极点''，因此可以在欧氏空间中
直接计算~\cite{Ji:2001wha}。

最近本文作者指出，非微扰 QCD 中的部分子物理属于本质上并非闵氏的一类问题~\cite{Ji:2013fga,Ji:2013dva,Hatta:2013gta,Xiong:2013bka,
Ji:2014gla,Ji:2014lra,Ji:2015jwa,Ji:2015qla,Xiong:2015nua}。这一观察基于两个要点：
1) 得益于渐近自由，光锥关联可以通过 QCD 微扰论用近光锥关联来近似；2) 洛伦兹助推是动力学的，
通过对强子态作助推，可以用它吸收场之间的动力学演化。于是所有剩余的动力学都包含在被助推的强子
波函数中，而不在探针算符里。这样，全部部分子物理都可以从大动量强子态中的等时关联获得，而这些
关联可以在欧氏蒙特卡罗模拟中计算。

该方法是一种有效场论，理由如下。对任意一个含时光锥算符 $O(\xi^-_1,\xi^-_2, ...)$，可以构造
无穷多个欧氏准算符 ${\cal O}_i(z_1,z_2,..)$，其中 $i$ 是所有此类算符的标记，$z_i$ 是场沿强子
动量方向（取为第三坐标方向）的空间坐标。这样的算符在大动量 $p_z$ 强子态中的矩阵元（例如在
格点 QCD 中计算）具有大动量展开~\cite{Ji:2014gla}
\begin{equation} \label{eq:fact}
\langle p|{\cal O}_i(z_1,z_2,..)|p\rangle
= C_{i0}(\alpha_s,p_z)\otimes\langle p|O(\xi^-_1,\xi^-_2, ...)|p\rangle + {\cal O}(1/p^2_z)\ ,
\end{equation}
其中我们略去了所有重正化标度，$C_{i0}$ 是硬系数。高阶修正可以系统地计算。由于在这个有效理论中
是用大动量去趋近光锥，故称为大动量有效理论，简称 LaMET。

值得指出的是，上面欧氏准算符的选取并不唯一。任何在无穷大动量极限下给出相同光前物理的准算符都
可用来探测所需的光锥算符，所有这些准算符构成一个普适类。实践中，普适类的存在允许人们对准算符作
优化选择，使 ${\cal O}(1/p_z^2)$ 修正尽可能小。这样一个例子是核子内总胶子极化的计算~\cite{Hatta:2013gta}。

近来，文献中对这一理论表现出很大兴趣，其中不乏对其正确性或有效性的质疑。这里我们提供若干讨论，
旨在进一步阐述该方法。文献~\cite{Carlson:2017gpk} 对 LaMET 方法中的因子化提出了质疑。循着
文献~\cite{Briceno:2017cpo}，我们明确证明该问题源于解析延拓的操作。文献~\cite{Rossi:2017muf}
基于矩的分析提出了幂次发散问题。我们证明这个问题在非局部做法中并不存在，那里唯一的幂次发散来自
威尔逊线的自能。最后，文献~\cite{Radyushkin:2017cyf,Orginos:2017kos} 提出了一种提取部分子分布的
新方法。我们证明 Ioffe 时间分布是从同一格点观测量提取部分子分布的另一种方式，因而与 LaMET 方法
完全等价。此外，在该方法中观察到的表观标度行为对收敛问题是否有帮助尚无定论。
